\BourbakiExerciseGroup

\begin{enumerate}
\def\labelenumi{\arabic{enumi})}
\item
  When M is a subgroup of the additive group R of real numbers, generalize the Euclidean division of polynomials in one indeterminate to the group algebra of M.
\item
  Let \(f\) be a polynomial \(\neq 0\) in \(\mathbf{A}[\mathbf{X}]\) of degree \(n\) and leading coefficient \(\alpha_0\) . Let \(\mathbf{M}\) be the submodule of \(\mathbf{A}[\mathbf{X}]\) (qua A-module) consisting of the polynomials in which the coefficient of the term of degree \(m\) (for arbitrary \(m\) in \(\mathbf{N}\) ) is divisible by \(\alpha_0^\mu\) , where \(\mu = \max(m - n + 1, 0)\) ; show that for any polynomial \(g\) in \(\mathbf{M}\) there exist two polynomials \(u, v\) in \(\mathbf{A}[\mathbf{X}]\) such that \(\deg v < n\) and \(g = uf + v\) .
\item
  \begin{enumerate}
  \def\labelenumii{\alph{enumii})}
  \tightlist
  \item
    In the algebra A{[}X{]} of polynomials in one indeterminate over A, the mapping \((u,v)\mapsto u(v)\) is an internal law of composition; show that this law is associative and distributive on the left with respect to addition and multiplication in A{[}X{]}.
  \end{enumerate}
\end{enumerate}

\begin{enumerate}
\def\labelenumi{\alph{enumi})}
\setcounter{enumi}{1}
\item
  If A is an integral domain show that the relation \(u(v)=0\) implies that u=0 or v is constant, and that when \(u\neq0\) and \(\deg v>0\) , the degree of \(u(v)\) is equal to the product of the degrees of u and v.
\item
  From now on assume that A is a field. Let u and v be two polynomials of degree \textgreater0 in A{[}X{]} and f a polynomial of degree \textgreater0; show that if q is the quotient and r the remainder in the Euclidean division of u by v, \(q(f)\) and \(r(f)\) are the quotient and remainder in the Euclidean division of \(u(f)\) by \(v(f)\) .
\item
  For each polynomial \(f\) of \(\mathbf{A}[\mathbf{X}]\) denote by \(\mathrm{I}(f)\) the set of all polynomials of the form \(u(f)\) , where \(u\) ranges over \(\mathbf{A}[\mathbf{X}]\) ; this is a subring of \(\mathbf{A}[\mathbf{X}]\) . For \(\mathrm{I}(f) = \mathrm{I}(g)\) it is necessary and sufficient that \(g = \lambda f + \mu\) , where \(\lambda \neq 0\) and \(\mu\) are in \(\mathbf{A}\) .
\item
  Let \(f\) and \(g\) be two polynomials of degree \(> 0\) in \(\mathbf{A}[\mathbf{X}]\) ; show that \(I(f) \cap I(g)\) is either equal to \(\mathbf{A}\) or of the form \(I(h)\) , where \(h\) is a polynomial of degree \(> 0\) (excluding the first possibility, consider in \(I(f) \cap I(g)\) a polynomial \(h\) of least possible degree \(> 0\) ).
\end{enumerate}

\begin{enumerate}
\def\labelenumi{\arabic{enumi})}
\setcounter{enumi}{3}
\tightlist
\item
  \begin{enumerate}
  \def\labelenumii{\alph{enumii})}
  \tightlist
  \item
    Suppose that for every non-zero \(n \in \mathbf{Z}\) , the relation \(n\xi = 0\) implies \(\xi = 0\) in \(\mathbf{A}\) . Show that the mapping \(f \mapsto f(\mathrm{D}_1, \ldots, \mathrm{D}_n)\) is an injective homomorphism of the polynomial algebra \(\mathbf{A}[\mathbf{X}_1, \ldots, \mathbf{X}_n] = \mathbf{E}\) into the algebra (over \(\mathbf{A}\) ) of endomorphisms of the \(\mathbf{A}\) -module \(\mathbf{E}\) (use induction on \(n\) ).
  \end{enumerate}
\end{enumerate}

\begin{enumerate}
\def\labelenumi{\alph{enumi})}
\setcounter{enumi}{1}
\tightlist
\item
  Let \(m\) be an integer \(\geqslant 1\) such that \(m: 1 = 0\) in \(\mathbf{A}\) . Show that we have \(\mathrm{D}_i^m = 0\) for every partial derivation \(\mathrm{D}_i\) ( \(1 \leqslant i \leqslant n\) ) in \(\mathbf{A}[\mathbf{X}_1, \ldots, \mathbf{X}_n]\) .
\end{enumerate}

\begin{enumerate}
\def\labelenumi{\arabic{enumi})}
\setcounter{enumi}{4}
\tightlist
\item
  Let \(u \in \mathbf{A}[(\mathbf{X}_i)_{i \in I}]\) and \(r\) an integer \(\geqslant 0\) , and suppose that in \(\mathbf{A}\) for each \(n > 0\) , the relation \(n\xi = 0\) implies \(\xi = 0\) . If we have the relation
\end{enumerate}

\[
\sum_ {i \in I} \left(\mathrm{D} _ {i} u\right) (\mathbf {X}) \mathbf {X} _ {i} = r u (\mathbf {X}),
\]

then \(u\) is homogeneous of degree \(r\) .

\begin{enumerate}
\def\labelenumi{\arabic{enumi})}
\setcounter{enumi}{5}
\tightlist
\item
  Let K be a commutative field, I a set and L the free associative algebra on the set I over the ring K (III, p.~448, Def. 2) with the graduation defined in III, p.~458. We define the degree of an element of L and the notion of homogeneous element as in IV, p.~2. a) Show that \(\deg(uv) = \deg(u) + \deg(v)\) and \(\deg(u + v) \leqslant \sup(\deg(u), \deg(v))\) for \(u, v \in L\) .
\end{enumerate}

\begin{enumerate}
\def\labelenumi{\alph{enumi})}
\setcounter{enumi}{1}
\item
  Given \(u, v, q, q' \in L\) such that \(\deg(u) \geqslant \deg(v)\) and \(\deg(qu - q'v) < \deg(ku)\) ; show that there exists \(q'' \in L\) such that \(\deg(u - q''v) < \deg(u)\) .
\item
  Given \(u, v, q, q' \in L\) with \(\deg u \geqslant \deg v > \deg (qu - q'v)\) ; show that there exists \(q'' \in L\) such that \(\deg (u - q''v) < \deg v\) .
\item
  Given \(u, v, u', v' \in L\) , homogeneous and such that \(uv = u'v'\) and \(\deg v \geqslant \deg v' \geqslant 0\) ; there exists \(t \in L\) such that \(u' = ut, v = tv'\) .
\item
  Given \(u, v \in L\) , homogeneous, such that \(uv = vu\) , there exists \(w \in L\) , \(a, b \in K\) and \(m, n \in K\) such that \(u = aw^m\) , \(v = bw^n\) .
\item
  Let \(u, v \in L\) be such that there exist \(a, b \in L\) , non-zero, with \(au = bv\) ; there exist \(m, d \in L\) such that the left multiples (resp. right divisors) common to \(u\) and \(v\) are left multiples (resp. right divisors) of \(m\) (resp. \(d\) ).
\end{enumerate}

